%% ============================================================
\chapter{Міжмолекулярні взаємодії}
%% ============================================================

Міжмолекулярні взаємодії --- фундаментальна складова фізико-хімічних процесів, що визначає властивості конденсованих фаз, стабільність комплексів, адсорбцію, а також біомолекулярне впізнавання. Розрахунок енергії взаємодії між двома фрагментами системи є одним з основних тестів для будь-якого квантово-хімічного методу.

%% --------------------------------------------------------
\section{Фізична природа міжмолекулярних взаємодій}
%% --------------------------------------------------------

Міжмолекулярні сили за своєю природою набагато слабші за хімічні (ковалентні) зв'язки. Типові енергії складають від $0.1$~кДж/моль для слабких дисперсійних взаємодій до 40~кДж/моль для міцних водневих зв'язків, тоді як енергія ковалентного зв'язку сягає 400~кДж/моль і більше.

\subsection{Класифікація міжмолекулярних взаємодій}

Міжмолекулярні взаємодії можна класифікувати за кількома критеріями:

\paragraph{За фізичною природою:}
\begin{itemize}
    \item \textbf{Електростатичні} --- взаємодія між постійними мультиполями (диполь-диполь, іон-диполь, квадруполь-квадруполь)
    \item \textbf{Індукційні} --- взаємодія між постійним мультиполем і індукованим
    \item \textbf{Дисперсійні} --- взаємодія між миттєвими (індукованими) диполями
    \item \textbf{Обмінні} --- відштовхування, спричинене принципом Паулі
    \item \textbf{Переносу заряду} --- делокалізація електронної густини між фрагментами
\end{itemize}

\paragraph{За відстанню:}
Електростатичні та індукційні внески спадають як \(R^{-n}\) (де \(n = 1, 2, 3, \ldots\) залежить від мультипольності), дисперсійні --- як \(R^{-6}\), обмінні --- експоненційно.

\paragraph{За специфічністю:}
\begin{itemize}
    \item \textbf{Неспецифічні} --- ван-дер-ваальсові сили (завжди присутні)
    \item \textbf{Специфічні} --- водневі зв'язки, галогенні зв'язки, \(\pi\)-\(\pi\) стекінг
\end{itemize}

\subsection{Водневий зв'язок}

Водневий зв'язок \( \mathrm{X{-}H \cdots Y} \) утворюється, коли атом водню, ковалентно зв'язаний з електронегативним атомом X (\ce{N}, \ce{O}, \ce{F}), взаємодіє з іншим електронегативним атомом Y, що має неподілену електронну пару.

Характерні ознаки водневого зв'язку:
\begin{itemize}
    \item Енергія: $5 - 40$~кДж/моль.
    \item Спрямованість: лінійна або майже лінійна геометрія.
    \item Відстань \( \mathrm{H \cdots Y} \): $1.5 - 2.5$~\AA (коротша за суму ван-дер-ваальсових радіусів).
    \item Подовження зв'язку \(\mathrm{X{-}H}\) та зсув його валентного коливання в червону область.
\end{itemize}

%% --------------------------------------------------------
\subsection{Дисперсійні взаємодії (сили Лондона)}
%% --------------------------------------------------------

Дисперсійні сили виникають через миттєві флуктуації електронної густини. Навіть у неполярних молекулах у кожен момент часу розподіл електронів не є ідеально симетричним, що створює миттєвий диполь. Цей миттєвий диполь індукує диполь у сусідній молекулі, і взаємодія між ними приводить до притягання.

Класичне наближення дає для двох атомів:
\[
E_{\text{disp}} = -\frac{C_6}{R^6},
\]
де \(C_6\) --- дисперсійний коефіцієнт, що залежить від поляризовності взаємодіючих частинок.

Дисперсійні сили:
\begin{itemize}
    \item Завжди присутні (універсальні).
    \item Не мають спрямованості.
    \item Швидко спадають з відстанню (\(\sim R^{-6}\)).
    \item Домінують у взаємодії неполярних молекул.
    \item Складають значну частину навіть у полярних системах.
\end{itemize}

%% --------------------------------------------------------
\section{Енергія взаємодії та суперпозиція базисних функцій}
%% --------------------------------------------------------

\subsection{Визначення енергії взаємодії}

Енергія взаємодії двох молекул \( A \) і \( B \) визначається як різниця між енергією димеру та сумою енергій ізольованих фрагментів:
\[
E_{\text{int}} = E_{AB} - (E_A + E_B).
\]

Ця величина має бути:
\begin{itemize}
    \item \textbf{Від'ємною} для стабілізуючої взаємодії (притягання).
    \item \textbf{Додатною} для дестабілізуючої взаємодії (відштовхування).
    \item \textbf{Незалежною від базису} (у ідеалі).
\end{itemize}

\subsection{Проблема перекриття базисів (BSSE)}

При використанні скінченного базисного набору виникає так зване \textbf{перекриття базисів} (Basis Set Superposition Error, BSSE). Суть проблеми полягає в наступному:

\begin{enumerate}
    \item При розрахунку димеру AB кожен фрагмент може використовувати базисні функції не лише <<свої>>, а й <<позичені>> від іншого фрагмента.
    \item Це штучно збільшує ефективний базисний набір для кожного фрагмента.
    \item Енергія димеру штучно занижується (стабілізується).
    \item При розрахунку ізольованих фрагментів кожен має доступ лише до власних базисних функцій.
    \item Різниця \(E_{AB} - (E_A + E_B)\) стає завищеною за абсолютною величиною.
\end{enumerate}

\textbf{Результат:} Енергія взаємодії переоцінюється (здається сильнішою, ніж насправді), особливо для малих базисних наборів.

\subsection{Метод контрпоїзу (Counterpoise Correction)}

Для усунення BSSE застосовують \textbf{корекцію за методом контрпоїзу}, яку вперше запропонували Бойс і Бернард (Boys and Bernardi, 1970).

\paragraph{Ідея методу:}
Розраховувати енергії ізольованих фрагментів не у власному базисі, а в <<димерному>> базисі, що включає базисні функції обох фрагментів:

\[
E_{\text{int}}^{\text{CP}} = E_{AB}(AB) - \left[ E_A(AB) + E_B(AB) \right],
\]

де:
\begin{itemize}
    \item \(E_{AB}(AB)\) --- енергія димеру в базисі димеру.
    \item \(E_A(AB)\) --- енергія фрагмента A в базисі димеру (позиції B залишаються як «ghost atoms»).
    \item \(E_B(AB)\) --- енергія фрагмента B в базісі димеру (позиції A як «ghost atoms»).
\end{itemize}

\textbf{Ghost atoms} --- це віртуальні центри, на яких розміщені базисні функції, але відсутні ядра та електрони.

\paragraph{Величина BSSE:}
\[
\text{BSSE} = E_{\text{int}} - E_{\text{int}}^{\text{CP}} = \left[E_A - E_A(AB)\right] + \left[E_B - E_B(AB)\right].
\]

Типові значення BSSE:
\begin{itemize}
    \item Базис STO-3G: 10--50\% від \(E_{\text{int}}\)
    \item Базис 6-31G*: 5--20\%
    \item Базис aug-cc-pVTZ: 1--5\%
    \item Базис aug-cc-pVQZ: < 1\%
\end{itemize}

\paragraph{Зауваження:}
\begin{itemize}
    \item Корекція CP завжди зменшує (за абсолютною величиною) енергію взаємодії
    \item Для повних базисів BSSE $\to$ 0 і корекція не потрібна
    \item Корекція CP не ідеальна --- може <<перекоригувати>> для дуже малих базисів
    \item Рекомендується використовувати достатньо великі базиси з дифузними функціями
\end{itemize}

%% --------------------------------------------------------
\section{Практичний приклад: димер води з корекцією BSSE}
%% --------------------------------------------------------

Розглянемо водневий димер \(\mathrm{(H_2O)_2}\). Ця система є класичним прикладом водневого зв'язку, у якому одночасно проявляються електростатичні, індукційні, обмінні та дисперсійні внески.

\subsection{Геометрія димеру води}

Стабільна конфігурація має наступну структуру:
\begin{itemize}
    \item Одна молекула води (донор) орієнтована так, що один з атомів \ce{H} спрямований до атома \ce{O} іншої молекули (акцептора).
    \item Водневий зв'язок має майже лінійну геометрію O--H···O.
    \item Відстань O···O становить близько $2.9$~\AA.
    \item Відстань H···O становить близько $1.9$~\AA.
\end{itemize}

\subsection{Код для розрахунку з корекцією BSSE}

\begin{minted}{python}
#!/usr/bin/env python3
"""
Розрахунок енергії взаємодії димеру води з корекцією BSSE
методом контрпоїзу (counterpoise correction)
"""

import numpy as np
from pyscf import gto, scf

# Геометрія димеру води (в ангстремах)
# Молекула A: донор протону
# Молекула B: акцептор протону
mol_dimer = gto.M(
    atom='''
    O  0.000000  0.000000  0.000000
    H  0.758602  0.000000  0.504284
    H  0.758602  0.000000 -0.504284
    O  2.900000  0.000000  0.000000
    H  3.300000  0.758602  0.000000
    H  3.300000 -0.758602  0.000000
    ''',
    basis='aug-cc-pVDZ',
    verbose=3
)

# Молекула A (перші 3 атоми)
mol_A = gto.M(
    atom='''
    O  0.000000  0.000000  0.000000
    H  0.758602  0.000000  0.504284
    H  0.758602  0.000000 -0.504284
    ''',
    basis='aug-cc-pVDZ',
    verbose=3
)

# Молекула B (останні 3 атоми)
mol_B = gto.M(
    atom='''
    O  2.900000  0.000000  0.000000
    H  3.300000  0.758602  0.000000
    H  3.300000 -0.758602  0.000000
    ''',
    basis='aug-cc-pVDZ',
    verbose=3
)

# Молекула A з ghost-атомами B (для CP-корекції)
mol_A_ghost = gto.M(
    atom='''
    O  0.000000  0.000000  0.000000
    H  0.758602  0.000000  0.504284
    H  0.758602  0.000000 -0.504284
    ghost:O  2.900000  0.000000  0.000000
    ghost:H  3.300000  0.758602  0.000000
    ghost:H  3.300000 -0.758602  0.000000
    ''',
    basis='aug-cc-pVDZ',
    verbose=3
)

# Молекула B з ghost-атомами A
mol_B_ghost = gto.M(
    atom='''
    ghost:O  0.000000  0.000000  0.000000
    ghost:H  0.758602  0.000000  0.504284
    ghost:H  0.758602  0.000000 -0.504284
    O  2.900000  0.000000  0.000000
    H  3.300000  0.758602  0.000000
    H  3.300000 -0.758602  0.000000
    ''',
    basis='aug-cc-pVDZ',
    verbose=3
)

# HF розрахунки
print("="*60)
print("Розрахунок енергій на рівні Hartree-Fock")
print("="*60)

mf_dimer = scf.RHF(mol_dimer).run()
E_dimer = mf_dimer.e_tot

mf_A = scf.RHF(mol_A).run()
E_A = mf_A.e_tot

mf_B = scf.RHF(mol_B).run()
E_B = mf_B.e_tot

mf_A_ghost = scf.RHF(mol_A_ghost).run()
E_A_ghost = mf_A_ghost.e_tot

mf_B_ghost = scf.RHF(mol_B_ghost).run()
E_B_ghost = mf_B_ghost.e_tot

# Розрахунок енергій взаємодії
E_int_uncorrected = (E_dimer - E_A - E_B) * 627.509  # ккал/моль
E_int_CP = (E_dimer - E_A_ghost - E_B_ghost) * 627.509
BSSE = E_int_uncorrected - E_int_CP

print("\n" + "="*60)
print("РЕЗУЛЬТАТИ (Hartree-Fock/aug-cc-pVDZ)")
print("="*60)
print(f"E(димер)              = {E_dimer:.8f} Hartree")
print(f"E(A)                  = {E_A:.8f} Hartree")
print(f"E(B)                  = {E_B:.8f} Hartree")
print(f"E(A з ghost B)        = {E_A_ghost:.8f} Hartree")
print(f"E(B з ghost A)        = {E_B_ghost:.8f} Hartree")
print(f"\nE_int (некориговане)  = {E_int_uncorrected:.2f} ккал/моль")
print(f"E_int (з CP-корекцією) = {E_int_CP:.2f} ккал/моль")
print(f"BSSE                  = {BSSE:.2f} ккал/моль")
print(f"BSSE (%)              = {abs(BSSE/E_int_uncorrected)*100:.1f}%")

# Порівняння з MP2
print("\n" + "="*60)
print("Порівняння з MP2")
print("="*60)

from pyscf import mp

mp2_dimer = mp.MP2(mf_dimer).run()
E_dimer_mp2 = mp2_dimer.e_tot

mp2_A = mp.MP2(mf_A).run()
E_A_mp2 = mp2_A.e_tot

mp2_B = mp.MP2(mf_B).run()
E_B_mp2 = mp2_B.e_tot

mp2_A_ghost = mp.MP2(mf_A_ghost).run()
E_A_ghost_mp2 = mp2_A_ghost.e_tot

mp2_B_ghost = mp.MP2(mf_B_ghost).run()
E_B_ghost_mp2 = mp2_B_ghost.e_tot

E_int_mp2_uncorr = (E_dimer_mp2 - E_A_mp2 - E_B_mp2) * 627.509
E_int_mp2_CP = (E_dimer_mp2 - E_A_ghost_mp2 - E_B_ghost_mp2) * 627.509
BSSE_mp2 = E_int_mp2_uncorr - E_int_mp2_CP

print(f"E_int MP2 (некориговане)   = {E_int_mp2_uncorr:.2f} ккал/моль")
print(f"E_int MP2 (з CP-корекцією) = {E_int_mp2_CP:.2f} ккал/моль")
print(f"BSSE MP2                   = {BSSE_mp2:.2f} ккал/моль")
print(f"BSSE MP2 (%)               = {abs(BSSE_mp2/E_int_mp2_uncorr)*100:.1f}%")

print("\n" + "="*60)
print("Експериментальне значення: ~5 ккал/моль (21 кДж/моль)")
print("="*60)
\end{minted}

\subsection{Очікувані результати}

Для димеру води очікуються наступні значення:

\begin{itemize}
    \item \textbf{HF/aug-cc-pVDZ:}
    \begin{itemize}
        \item \(E_{\text{int}}\) без корекції: -6.5 ккал/моль
        \item \(E_{\text{int}}\) з CP: -4.8 ккал/моль
        \item BSSE: 1.7 ккал/моль (26\%)
    \end{itemize}

    \item \textbf{MP2/aug-cc-pVDZ:}
    \begin{itemize}
        \item \(E_{\text{int}}\) без корекції: -5.4 ккал/моль
        \item \(E_{\text{int}}\) з CP: -5.0 ккал/моль
        \item BSSE: 0.4 ккал/моль (7\%)
    \end{itemize}
\end{itemize}

\textbf{Висновки:}
\begin{enumerate}
    \item HF недооцінює енергію взаємодії через відсутність дисперсії
    \item MP2 дає значно кращі результати, близькі до експерименту
    \item BSSE для aug-cc-pVDZ помірна, але все ще значуща
    \item Для кількісних результатів потрібні більші базиси (aug-cc-pVTZ або більше)
\end{enumerate}

%% --------------------------------------------------------
\section{Міжмолекулярні взаємодії в різних методах}
%% --------------------------------------------------------

\subsection{Метод Гартрі---Фока}

Метод \texttt{HF} описує лише обмінну взаємодію та електростатику на рівні середнього поля, але не враховує кореляційної енергії, тому він \textbf{не здатен відтворювати дисперсійні сили}.

\paragraph{Що враховує HF:}
\begin{itemize}
    \item Електростатична взаємодія між розподілами заряду
    \item Обмінне відштовхування (принцип Паулі)
    \item Частково індукційну взаємодію
\end{itemize}

\paragraph{Що не враховує HF:}
\begin{itemize}
    \item Дисперсійні сили (взагалі)
    \item Точну індукцію (лише в наближенні середнього поля)
    \item Кореляцію електронів у фрагментах
\end{itemize}

\paragraph{Наслідки:}
\begin{itemize}
    \item Енергія взаємодії неполярних систем (наприклад, димер метану) майже нульова або додатна
    \item Недооцінка стабільності всіх слабко зв'язаних систем
    \item Помилки зростають зі збільшенням розміру системи
\end{itemize}

\subsection{Пост-HF методи}

Методи \texttt{MP2}, \texttt{CCSD}, \texttt{CCSD(T)} та інші пост-HF підходи враховують електронну кореляцію і тому природно відтворюють дисперсійні взаємодії.

\paragraph{MP2 (Møller--Plesset 2-го порядку):}
\begin{itemize}
    \item Найпростіший кореляційний метод
    \item Масштабується як \(O(N^5)\)
    \item Як правило, \textbf{переоцінює} дисперсійні взаємодії (на 10--20\%)
    \item Добре працює для систем з водневими зв'язками
    \item Може давати помилки для \(\pi\)-\(\pi\) стекінгу
\end{itemize}

\paragraph{CCSD(T) (Coupled Cluster з синглетами, дублетами і пертурбативними триплетами):}
\begin{itemize}
    \item <<Золотий стандарт>> квантової хімії
    \item Масштабується як \(O(N^7)\)
    \item Точність: 0.1--0.3 ккал/моль для невеликих систем
    \item Еталонний метод для калібрування інших підходів
    \item Обмежений системами до 10--20 важких атомів
\end{itemize}

\subsection{DFT і дисперсійні поправки}

Стандартні функціонали DFT (особливо локальні та напів-локальні) \textbf{не містять дисперсійних взаємодій} за побудовою, оскільки дисперсія --- це нелокальний кореляційний ефект.

\paragraph{Проблеми стандартних функціоналів:}
\begin{itemize}
    \item \texttt{LDA}, \texttt{PBE}: майже повна відсутність дисперсії
    \item \texttt{B3LYP}: частково враховує дисперсію через HF-обмін, але недостатньо
    \item \texttt{M06-2X}, \texttt{ωB97X}: краща поведінка через параметризацію на недисперсійних системах
\end{itemize}

\paragraph{Емпіричні дисперсійні поправки:}

\textbf{DFT-D3 (Grimme, 2010):}
\[
E_{\text{DFT-D3}} = E_{\text{DFT}} + E_{\text{disp}},
\]
де
\[
E_{\text{disp}} = -\sum_{n=6,8,10} s_n \sum_{i<j} \frac{C_n^{ij}}{R_{ij}^n} f_{\text{damp}}(R_{ij}).
\]

Переваги:
\begin{itemize}
    \item Малі обчислювальні витрати (майже безкоштовно)
    \item Добре параметризовані коефіцієнти
    \item Універсальність для різних функціоналів
\end{itemize}

\textbf{DFT-D4 (Grimme, 2019):}
\begin{itemize}
    \item Врахування заряду та геометрії
    \item Краща точність для заряджених систем та металів
\end{itemize}

\textbf{VV10 (нелокальний кореляційний функціонал):}
\[
E_c^{\text{nl}} = \int \rho(\mathbf{r}) \Phi[\rho(\mathbf{r}), \nabla\rho(\mathbf{r}), \rho(\mathbf{r}')] \, d\mathbf{r},
\]
де \(\Phi\) --- нелокальне ядро.

Сучасні підходи, такі як \texttt{ωB97X-D}, \texttt{ωB97M-V}, або \texttt{rVV10}, включають ці поправки безпосередньо у формалізм функціоналу.

\subsection{Практичне порівняння методів}

\begin{minted}{python}
#!/usr/bin/env python3
"""
Порівняння різних методів для розрахунку
енергії взаємодії димеру метану (неполярна система)
"""

from pyscf import gto, scf, mp, dft

# Геометрія димеру метану
mol_dimer = gto.M(
    atom='''
    C  0.0  0.0  0.0
    H  0.629118  0.629118  0.629118
    H -0.629118 -0.629118  0.629118
    H -0.629118  0.629118 -0.629118
    H  0.629118 -0.629118 -0.629118
    C  3.5  0.0  0.0
    H  4.129118  0.629118  0.629118
    H  2.870882 -0.629118  0.629118
    H  2.870882  0.629118 -0.629118
    H  4.129118 -0.629118 -0.629118
    ''',
    basis='aug-cc-pVDZ'
)

mol_monomer = gto.M(
    atom='''
    C  0.0  0.0  0.0
    H  0.629118  0.629118  0.629118
    H -0.629118 -0.629118  0.629118
    H -0.629118  0.629118 -0.629118
    H  0.629118 -0.629118 -0.629118
    ''',
    basis='aug-cc-pVDZ'
)

def compute_interaction_energy(method_name, mf_dimer, mf_monomer):
    """Розрахунок енергії взаємодії"""
    E_dimer = mf_dimer.kernel()
    E_monomer = mf_monomer.kernel()
    E_int = (E_dimer - 2*E_monomer) * 627.509  # ккал/моль
    print(f"{method_name:20s}: E_int = {E_int:8.3f} ккал/моль")
    return E_int

print("="*60)
print("Енергія взаємодії димеру метану (aug-cc-pVDZ)")
print("="*60)

# Hartree-Fock
mf_hf_d = scf.RHF(mol_dimer)
mf_hf_m = scf.RHF(mol_monomer)
compute_interaction_energy("HF", mf_hf_d, mf_hf_m)

# MP2
mp2_d = mp.MP2(scf.RHF(mol_dimer).run())
mp2_m = mp.MP2(scf.RHF(mol_monomer).run())
compute_interaction_energy("MP2", mp2_d, mp2_m)

# B3LYP (без дисперсії)
mf_b3lyp_d = dft.RKS(mol_dimer)
mf_b3lyp_d.xc = 'B3LYP'
mf_b3lyp_m = dft.RKS(mol_monomer)
mf_b3lyp_m.xc = 'B3LYP'
compute_interaction_energy("B3LYP", mf_b3lyp_d, mf_b3lyp_m)

# PBE (без дисперсії)
mf_pbe_d = dft.RKS(mol_dimer)
mf_pbe_d.xc = 'PBE'
mf_pbe_m = dft.RKS(mol_monomer)
mf_pbe_m.xc = 'PBE'
compute_interaction_energy("PBE", mf_pbe_d, mf_pbe_m)

# M06-2X (частково враховує дисперсію)
mf_m06_d = dft.RKS(mol_dimer)
mf_m06_d.xc = 'M062X'
mf_m06_m = dft.RKS(mol_monomer)
mf_m06_m.xc = 'M062X'
compute_interaction_energy("M06-2X", mf_m06_d, mf_m06_m)

print("\nЕкспериментальне значення: ~0.5 ккал/моль")
print("Еталонне CCSD(T): ~0.52 ккал/моль")
\end{minted}

\subsection{Очікувані результати для димеру метану}

\begin{table}[h]
\centering
\begin{tabular}{lcc}
\hline
\textbf{Метод} & \textbf{E$_{\text{int}}$ (ккал/моль)} & \textbf{Коментар} \\
\hline
HF              & +0.2  & Відштовхування (немає дисперсії) \\
B3LYP           & +0.1  & Майже немає взаємодії \\
PBE             & +0.3  & Відштовхування \\
M06-2X          & -0.3  & Частково враховує дисперсію \\
MP2             & -0.6  & Переоцінка дисперсії \\
\hline
Експеримент     & -0.5  & Еталонне значення \\
CCSD(T)         & -0.52 & Золотий стандарт \\
\hline
\end{tabular}
\caption{Енергія взаємодії димеру метану різними методами}
\end{table}

%% --------------------------------------------------------
\section{Розкладання енергії взаємодії: SAPT-аналіз}
%% --------------------------------------------------------

Для глибшого розуміння природи міжмолекулярних зв'язків застосовують \textbf{метод симетризованої теорії збурень} (\textit{Symmetry-Adapted Perturbation Theory}, SAPT).

На відміну від підходу прямого віднімання енергій, SAPT дозволяє розкласти енергію взаємодії на фізично змістовні компоненти:

\[
E_{\text{int}} = E_{\text{elst}} + E_{\text{exch}} + E_{\text{ind}} + E_{\text{disp}} + \delta_{\text{HF}} + \cdots
\]

\subsection{Компоненти SAPT}

\paragraph{1. Електростатична енергія (\(E_{\text{elst}}\)):}
\begin{itemize}
    \item Взаємодія між незбуреними розподілами заряду фрагментів
    \item Може бути як притягальною (різнойменні заряди), так і відштовхувальною
    \item Залежить від мультипольних моментів
    \item Спадає як \(R^{-1}\) (монополь-монополь), \(R^{-2}\) (монополь-диполь), \(R^{-3}\) (диполь-диполь) тощо
\end{itemize}

\paragraph{2. Обмінна енергія (\(E_{\text{exch}}\)):}
\begin{itemize}
    \item Відштовхування внаслідок принципу Паулі
    \item Завжди додатна (дестабілізуюча)
    \item Виникає при перекритті хвильових функцій фрагментів
    \item Спадає експоненційно: \(\sim \exp(-\alpha R)\)
    \item Домінує на малих відстанях
\end{itemize}

\paragraph{3. Індукційна енергія (\(E_{\text{ind}}\)):}
\begin{itemize}
    \item Поляризація одного фрагмента електричним полем іншого
    \item Завжди притягальна (стабілізуюча)
    \item Включає ефекти переносу заряду
    \item Спадає як \(R^{-4}\) (диполь-індукований диполь)
    \item Важлива для систем з полярними групами
\end{itemize}

\paragraph{4. Дисперсійна енергія (\(E_{\text{disp}}\)):}
\begin{itemize}
    \item Кореляція миттєвих флуктуацій електронної густини
    \item Завжди притягальна
    \item Спадає як \(R^{-6}\)
    \item Присутня у всіх системах
    \item Домінує для неполярних молекул
    \item Зростає з поляризовністю
\end{itemize}

\paragraph{5. Корекція вищих порядків (\(\delta_{\text{HF}}\)):}
\begin{itemize}
    \item Різниця між повним HF і сумою електростатики + обміну + індукції
    \item Враховує зв'язок між внесками
    \item Зазвичай невелика (5--10\% від загальної енергії)
\end{itemize}

\subsection{Рівні теорії SAPT}

\paragraph{SAPT0:}
\begin{itemize}
    \item Найпростіший варіант
    \item Електростатика та обмін на рівні HF
    \item Індукція та дисперсія в найнижчому порядку
    \item Масштабування: \(O(N^5)\)
    \item Точність: ~20\% для слабких взаємодій
\end{itemize}

\paragraph{SAPT2:}
\begin{itemize}
    \item Включає корекції другого порядку
    \item Точність: ~5--10\%
    \item Масштабування: \(O(N^5)\)--\(O(N^6)\)
\end{itemize}

\paragraph{SAPT2+3:}
\begin{itemize}
    \item Додає внески третього порядку
    \item Точність: ~2--5\%
    \item Масштабування: \(O(N^7)\)
    \item Порівнянний з CCSD(T) за точністю
\end{itemize}

\subsection{SAPT у практиці розрахунків}

Розрахунки SAPT потребують спеціальних реалізацій. У PySCF існує базова підтримка SAPT-компонентів через модуль \texttt{pyscf.sapt}.

\begin{minted}{python}
#!/usr/bin/env python3
"""
SAPT-аналіз димеру води: розкладання енергії взаємодії
на фізичні компоненти
"""

import numpy as np
from pyscf import gto, scf, lib

# Примітка: PySCF не має повної реалізації SAPT
# Цей код демонструє концепцію та показує як обчислити
# деякі компоненти вручну

# Геометрія димеру води
atom_A = '''
O  0.000000  0.000000  0.000000
H  0.758602  0.000000  0.504284
H  0.758602  0.000000 -0.504284
'''

atom_B = '''
O  2.900000  0.000000  0.000000
H  3.300000  0.758602  0.000000
H  3.300000 -0.758602  0.000000
'''

basis = 'aug-cc-pVDZ'

# Створюємо молекули
mol_A = gto.M(atom=atom_A, basis=basis, verbose=0)
mol_B = gto.M(atom=atom_B, basis=basis, verbose=0)
mol_AB = gto.M(atom=atom_A + atom_B, basis=basis, verbose=0)

# HF розрахунки для окремих фрагментів
mf_A = scf.RHF(mol_A).run()
mf_B = scf.RHF(mol_B).run()
mf_AB = scf.RHF(mol_AB).run()

print("="*70)
print("SAPT-подібний аналіз димеру води")
print("="*70)

# 1. Повна енергія взаємодії на рівні HF
E_int_HF = (mf_AB.e_tot - mf_A.e_tot - mf_B.e_tot) * 627.509
print(f"\nПовна E_int (HF) = {E_int_HF:.3f} ккал/моль")

# 2. Електростатична компонента (наближення)
# Використовуємо густини окремих фрагментів
def compute_electrostatic(mol_A, mol_B, dm_A, dm_B):
    """
    Наближений розрахунок електростатичної енергії
    через взаємодію розподілів заряду
    """
    # Ядерно-електронна взаємодія A-B
    nuc_elec = 0.0
    coords_A = mol_A.atom_coords()
    coords_B = mol_B.atom_coords()
    charges_A = mol_A.atom_charges()
    charges_B = mol_B.atom_charges()

    # Ядро A - ядро B
    for i, (Za, Ra) in enumerate(zip(charges_A, coords_A)):
        for j, (Zb, Rb) in enumerate(zip(charges_B, coords_B)):
            R = np.linalg.norm(Ra - Rb)
            nuc_elec += Za * Zb / R

    # Тут має бути взаємодія електронних густин
    # (складна інтегральна процедура, спрощена для демонстрації)

    return nuc_elec * 627.509

E_elst_approx = compute_electrostatic(mol_A, mol_B,
                                      mf_A.make_rdm1(),
                                      mf_B.make_rdm1())

print(f"\n--- Наближена декомпозиція ---")
print(f"E_elst (наближення) ≈ {E_elst_approx:.3f} ккал/моль")

# 3. Для точного SAPT потрібні спеціалізовані модулі
print("\n" + "="*70)
print("Для повного SAPT-аналізу використовуйте:")
print("  - PSI4 (найкраща реалізація SAPT)")
print("  - MOLPRO")
print("  - SAPT2020 (окремий пакет)")
print("="*70)

# Демонстрація концепції: типові значення для димеру води
print("\n" + "="*70)
print("Типові SAPT0 результати для димеру води (aug-cc-pVDZ):")
print("="*70)
print(f"{'Компонента':<25s} {'Енергія (ккал/моль)':>20s}")
print("-"*70)
print(f"{'E_elst':<25s} {-9.2:>20.2f}")
print(f"{'E_exch':<25s} {+8.5:>20.2f}")
print(f"{'E_ind':<25s} {-2.1:>20.2f}")
print(f"{'E_disp':<25s} {-2.4:>20.2f}")
print(f"{'δ_HF':<25s} {-0.3:>20.2f}")
print("-"*70)
print(f"{'E_int (SAPT0)':<25s} {-5.5:>20.2f}")
print("-"*70)

# Графічне представлення внесків
print("\nВізуалізація внесків:")
print("  Притягання:")
print("    E_elst:  " + "█"*46 + f" {-9.2:.1f}")
print("    E_ind:   " + "█"*11 + f" {-2.1:.1f}")
print("    E_disp:  " + "█"*12 + f" {-2.4:.1f}")
print("  Відштовхування:")
print("    E_exch:  " + "█"*42 + f" {+8.5:.1f}")
print("  Сума:     " + "█"*28 + f" {-5.5:.1f}")

print("\n" + "="*70)
print("Фізична інтерпретація:")
print("="*70)
print("• Електростатика домінує у притяганні (~65%)")
print("• Обмін компенсує ~60% притягання")
print("• Індукція та дисперсія додають ~30% притягання")
print("• Результуюча взаємодія: помірний водневий зв'язок")
print("="*70)
\end{minted}

\subsection{Інтеграція з PSI4 для повного SAPT}

Для проведення повного SAPT-аналізу рекомендується використовувати PSI4, який має найкращу реалізацію SAPT. Нижче наведено приклад скрипта:

\begin{minted}{python}
#!/usr/bin/env python3
"""
Повний SAPT-аналіз димеру води за допомогою PSI4
(вимагає встановлення PSI4: conda install psi4 -c psi4)
"""

import psi4
import numpy as np

# Налаштування PSI4
psi4.set_memory('2 GB')
psi4.core.set_output_file('h2o_dimer_sapt.log', False)

# Геометрія димеру
dimer = psi4.geometry("""
0 1
O  0.000000  0.000000  0.000000
H  0.758602  0.000000  0.504284
H  0.758602  0.000000 -0.504284
--
0 1
O  2.900000  0.000000  0.000000
H  3.300000  0.758602  0.000000
H  3.300000 -0.758602  0.000000
units angstrom
""")

# Налаштування SAPT
psi4.set_options({
    'basis': 'aug-cc-pVDZ',
    'scf_type': 'df',           # Density fitting для швидкості
    'freeze_core': 'true',      # Заморожування кор-електронів
})

# Розрахунок SAPT0
print("="*70)
print("SAPT0 аналіз димеру води")
print("="*70)

energy_sapt = psi4.energy('sapt0', molecule=dimer)

# Отримання компонентів
elst = psi4.variable('SAPT ELST ENERGY') * 627.509
exch = psi4.variable('SAPT EXCH ENERGY') * 627.509
ind = psi4.variable('SAPT IND ENERGY') * 627.509
disp = psi4.variable('SAPT DISP ENERGY') * 627.509
total = psi4.variable('SAPT0 TOTAL ENERGY') * 627.509

print(f"\n{'Компонента':<25s} {'Енергія (ккал/моль)':>20s}")
print("-"*70)
print(f"{'Електростатика':<25s} {elst:>20.3f}")
print(f"{'Обмін':<25s} {exch:>20.3f}")
print(f"{'Індукція':<25s} {ind:>20.3f}")
print(f"{'Дисперсія':<25s} {disp:>20.3f}")
print("-"*70)
print(f"{'Повна енергія (SAPT0)':<25s} {total:>20.3f}")
print("="*70)

# Аналіз відносних внесків
total_attractive = abs(elst) + abs(ind) + abs(disp)
print(f"\nВідносні внески у притягання:")
print(f"  Електростатика: {abs(elst)/total_attractive*100:.1f}%")
print(f"  Індукція:       {abs(ind)/total_attractive*100:.1f}%")
print(f"  Дисперсія:      {abs(disp)/total_attractive*100:.1f}%")

# SAPT2 для кращої точності (опціонально)
print("\n" + "="*70)
print("SAPT2 аналіз (вища точність)")
print("="*70)

energy_sapt2 = psi4.energy('sapt2', molecule=dimer)
total_sapt2 = psi4.variable('SAPT2 TOTAL ENERGY') * 627.509

print(f"Повна енергія (SAPT2): {total_sapt2:.3f} ккал/моль")
print(f"Різниця SAPT2-SAPT0:   {total_sapt2-total:.3f} ккал/моль")
print("="*70)
\end{minted}

\subsection{Типові результати SAPT для різних систем}

\begin{table}[h]
\centering
\small
\begin{tabular}{lcccccc}
\hline
\textbf{Система} & \textbf{E$_{\text{elst}}$} & \textbf{E$_{\text{exch}}$} & \textbf{E$_{\text{ind}}$} & \textbf{E$_{\text{disp}}$} & \textbf{E$_{\text{int}}$} \\
\hline
(H$_2$O)$_2$  & $-9.2$  & $+8.5$  & $-2.1$  & $-2.4$  & $-5.5$ \\
(NH$_3$)$_2$  & $-6.8$  & $+6.2$  & $-1.5$  & $-1.8$  & $-4.1$ \\
(CH$_4$)$_2$  & $+0.1$  & $+0.8$  & $-0.1$  & $-1.2$  & $-0.5$ \\
Бензол димер  & $-1.8$  & $+5.2$  & $-0.6$  & $-5.1$  & $-2.8$ \\
(HF)$_2$      & $-12.5$ & $+10.8$ & $-2.8$  & $-1.9$  & $-6.8$ \\
\hline
\end{tabular}
\caption{SAPT0/aug-cc-pVDZ компоненти для різних димерів (ккал/моль)}
\end{table}

\paragraph{Спостереження:}
\begin{itemize}
    \item \textbf{Водневі зв'язки} (H$_2$O, NH$_3$, HF): електростатика домінує
    \item \textbf{Неполярні системи} (CH$_4$): дисперсія є єдиним джерелом притягання
    \item \textbf{$\pi$-$\pi$ стекінг} (бензол): дисперсія та електростатика приблизно рівні
    \item \textbf{Обмін}: завжди відштовхувальний, компенсує 60--90\% притягання
\end{itemize}

%% --------------------------------------------------------
\section{Energy Decomposition Analysis (EDA)}
%% --------------------------------------------------------

Альтернативним підходом до SAPT є \textbf{Energy Decomposition Analysis} (EDA), що базується на варіаційному принципі замість теорії збурень.

\subsection{Схема Морокуми (Morokuma)}

Класична схема Морокуми розкладає енергію взаємодії на: